\documentclass{article}

\usepackage{hyperref}
\usepackage[margin=1in,letterpaper]{geometry}
\usepackage[n,advantage,operators,sets,adversary,landau,probability,notions,logic,ff,mm,primitives,events,complexity,oracles,asymptotics,keys]{cryptocode}
\usepackage{amsfonts,amsmath,amssymb,amsthm}
\usepackage{algorithm}
\usepackage[noend]{algpseudocode}
\usepackage{array}
\usepackage{algpseudocodex}

\usepackage{multicol}
\usepackage{todonotes}
\usepackage{cleveref}
\usepackage{xspace}

\newcommand{\INDCCA}{\indcca}
\newcommand{\PRF}{\prf}

\newcommand{\keygen}{\pcalgostyle{KeyGen}}
\newcommand{\encaps}{\pcalgostyle{Encaps}}
\newcommand{\decaps}{\pcalgostyle{Decaps}}

% \newcommand{\sk}{sk}
% \newcommand{\pk}{pk}

\newcommand{\ct}{\pckeystyle{ct}}
\newcommand{\kemkeyspace}{\mathcal{K}}
\newcommand{\pkemsgspace}{\mathcal{M}}
\newcommand{\kem}{\pcalgostyle{KEM}}
\newcommand{\KEM}{\pcalgostyle{KEM}}
\newcommand{\REAL}{\pcnotionstyle{Real}}
\newcommand{\RANDOM}{\pcnotionstyle{Random}}
\newcommand{\PKE}{\pcalgostyle{PKE}}
\newcommand{\LEFT}{\pcnotionstyle{Left}}
\newcommand{\RIGHT}{\pcnotionstyle{Right}}

\newcommand{\nominalgroup}{\mathcal{N}}
\newcommand{\nominalgroupexpanded}{(\nelem, \ngen, p, \honestexponent , \exponent , \nexp)}
\newcommand{\honestexponent}{\varepsilon_h}
\newcommand{\exponent}{\varepsilon_u}
\newcommand{\nexp}{\textsf{exp}\xspace}
\newcommand{\nelem}{\mathcal{G}}
\newcommand{\ngen}{g}
\newcommand{\prfqueries}{q_e}
\newcommand{\ngalg}{\mathsf{NG}}

\newcommand{\sdh}{\text{SDH}\xspace}

\newcommand{\assign}{\leftarrow}

\newcommand{\hqb}{\pcalgostyle{HQB}}
\newcommand{\HQB}{\pcalgostyle{HQB}}

\newcommand{\QSH}{\pcalgostyle{QSH}}
\newcommand{\domain}{\textit{Label}}
\newcommand{\pr}{\pcnotionstyle{PR}}

\newcommand{\algor}{\textbf{Algorithm}\xspace}
\newcommand{\getsr}{{\:\leftarrow\hspace*{-3pt}\raisebox{.5pt}{$\scriptscriptstyle\$$}\:}}
\newcommand{\tor}{{\:\raisebox{.5pt}{$\scriptscriptstyle\$$}\hspace*{-3pt}\rightarrow\:}}
\newcommand{\foutput}{\rightarrow}
\newcommand{\ifthen}[1]{\textbf{if } #1 \textbf{ then}}
    \newcommand{\compare}{=}

    \usepackage[dvipsnames,table]{xcolor}
    \usepackage{adjustbox}
    \newcommand{\solidbox}[1]{\adjustbox{fbox}{\strut #1}}
    \newcommand{\graybox}[1]{\adjustbox{cframe=black!15, bgcolor=black!15}{\strut #1}}
    \newcommand{\highlightbox}[2][RoyalBlue!20]{\adjustbox{cframe=#1, bgcolor=#1}{\strut #2}}

    \theoremstyle{theorem}
    \newtheorem{theorem}{Theorem}
    \newtheorem{claim}{Claim}
    \theoremstyle{definition}
    \newtheorem{definition}{Definition}


    \newenvironment{gamefigure}[2]{
        \gdef\storeandloadcaption{#1}
        \gdef\storeandloadlabel{#2}
        \begin{figure}[ht!]
            \hrule
            \vspace{1mm}
        }{
            \hrule
            \caption{\storeandloadcaption} \label{\storeandloadlabel}
        \end{figure}
    }

    \newenvironment{games}[1]{
        \begin{multicols}{#1}
        }{
        \end{multicols}
    }

    \newenvironment{gameenv}[1]{
        \begin{algorithmic}
            \Statex \underline{#1}
            % \vspace{1mm}
        }{
        \end{algorithmic}
    }

    \newcommand{\gamebreak}{\columnbreak}

    \newcommand{\emptygame}{
        \begin{gameenv}{\null}
        \end{gameenv}
    }

    \title{StarFortress—  Hybrid Post-Quantum KEMs From SDH and IND-CCA}
    \author{Deirdre Connolly, Paul Grubbs}
    \date{\today}

    \begin{document}

    \maketitle

    \section{Introduction}

    \paragraph{Structure of this paper.}
    Notation and definitions are introduced \autoref{sec:prelims}.
    The HQB framework is introduced in \autoref{sec:hqb}; our generic
    construction -- named Heavy Quantum Bomber -- is a general framework
    that can be instantiated with different component algorithms.
    In \autoref{sec:hqb-proofs} we prove the IND-CCA security when the different
    components fail in different ways.
    Then in \autoref{sec:concrete} we give concrete instances of $\HQB$
    and show that they are secure.

    \section{Definitions}\label{sec:prelims}

    \subsection{Key encapsulation mechanisms}\label{sec:kems}

    \begin{definition}[KEM]\label{def:kem}
        A \emph{key encapsulation mechanism} scheme consists of three algorithms and set $\kemkeyspace$:
        \begin{itemize}
        \item $\keygen() \tor (\pk, \sk)$
        \item $\encaps(\pk) \tor (\ct, \key)$
        \item $\decaps(\sk, \ct) \to \key$
        \end{itemize}
    \end{definition}

    \begin{definition}[Correctness]\label{sec:kem-correctness}
        The correctness of a $\KEM$ imposes that, except with small probability drawn
        over the random coin space of $\KEM.\keygen$ and $\KEM.\encaps$, we have that
        $\KEM.\decaps$ correctly recovers the shared key produced by $\KEM.\encaps$.

        Formally, we say that a KEM is $\delta$-correct if:
        \[
            \Pr \left[\,  k \ne k' \, \left | \substack{(\sk, \pk) \sample \keygen(\,) \\ (k, c) \sample \encaps(\pk) \\ k' \gets \decaps(c, \sk) } \right. \right] \le \delta .
        \]
    \end{definition}

    \begin{definition}[$\indcca$ Security]\label{def:kem-security}
        Let $\KEM$ be a KEM and $\adv$ be an adversary. Define the advantage of $\adv$ in breaking the $\indcca$ security of $\KEM$ by
        \[ \advantage{\indcca}{\KEM}[(\adv)] = \Pr\left[ \indcca_\KEM^\adv.\REAL() \Rightarrow 1 \right] - \Pr\left[ \indcca_\KEM^\adv.\RANDOM() \Rightarrow 1 \right] \]
        \begin{pchstack}[center,space=1em]
            \begin{pcvstack}[boxed]
                \procedure[linenumbering]{$\indcca_\KEM^\adv.\REAL()$}{
                    (\pk, \sk) \getsr \KEM.\keygen() \\
                    (\ct^*, \highlightbox{$\key$}) \getsr \KEM.\encaps(\pk) \\
                    \pcreturn \adv^{\oracle}(\pk, \ct^*, \key)
                }
                \procedure[linenumbering]{$\oracle(\ct)$}{
                    \pcif \ct = \ct^* \pcthen \\
                    \pcind \pcreturn \bot \\
                    \pcelse \\
                    \pcind \pcreturn \KEM.\decaps(\sk, \ct)
                }
            \end{pcvstack}
            \begin{pcvstack}[boxed]
                \procedure[linenumbering]{$\indcca_\KEM^\adv.\RANDOM()$}{
                    (\pk, \sk) \getsr \KEM.\keygen() \\
                    (\ct^*, \_\_) \getsr \KEM.\encaps(\pk) \\
                    \highlightbox{$\key \getsr \KEM.\kemkeyspace$} \\
                    \pcreturn \adv^{\oracle}(\pk, \ct^*, \key)
                }
                \procedure[linenumbering]{$\oracle(\ct)$}{
                    \pcif \ct = \ct^* \pcthen \\
                    \pcind \pcreturn \bot \\
                    \pcelse \\
                    \pcind \pcreturn \KEM.\decaps(\sk, \ct)
                }
            \end{pcvstack}
        \end{pchstack}
    \end{definition}

    \subsection{Pseudorandom functions}\label{sec:prf}

    \begin{definition}[PRF]\label{def:prf}
        A \emph{PRF} is a function $F : \kemkeyspace \times \{0,1\}^* \to \kemkeyspace_o$.
    \end{definition}

    \begin{definition}[$\pr$ security for a PRF]\label{def:prf-security}
        Let $F$ be a PRF and $\adv$ be an adversary. Define the advantage of $\adv$ in breaking the $\pr$ (pseudorandomness) security of $F$ by
        \[ \advantage{\pr}{F}[(\adv)] = \Pr\left[ \pr_{F}^\adv.\REAL() \Rightarrow 1 \right] - \Pr\left[ \pr_{F}^\adv.\RANDOM() \Rightarrow 1 \right] \]

        \begin{pchstack}[center,space=1em]
            \begin{pcvstack}[boxed]
                \procedure[linenumbering]{$\pr_{F}^\adv.\REAL()$}{
                    X \gets \emptyset \\
                    \highlightbox{$k \getsr \kemkeyspace$} \\
                    \pcreturn \adv^{\oracle}
                }
                \procedure[linenumbering]{$\oracle(x)$}{
                    \pcif x \in X \pcthen \mbox{abort} \\
                    X \gets X \cup \{x\} \\
                    \highlightbox{$y \gets F(k, x)$} \\
                    \pcreturn y
                }
            \end{pcvstack}
            \begin{pcvstack}[boxed]
                \procedure[linenumbering]{$\pr_{F}^\adv.\RANDOM()$}{
                    X \gets \emptyset \\
                    \pcreturn \adv^{\oracle}
                }
                \procedure[linenumbering]{$\oracle(x)$}{
                    \pcif x \in X \pcthen \mbox{abort} \\
                    X \gets X \cup \{x\} \\
                    \highlightbox{$y \getsr \kemkeyspace_o$} \\
                    \pcreturn y
                }
            \end{pcvstack}
        \end{pchstack}
    \end{definition}

    \todo[inline]{Do we need split-key PRFs?}

    \subsection{Nominal Group}

    In this section, we recall the definition of a nominal group, as of \cite{EC:ABHKLR21}.

    \begin{definition}[Nominal Group \cite{EC:ABHKLR21}]
        A nominal group $\nominalgroup =\nominalgroupexpanded$ consists of an efficiently recognizable finite set of elements $\mathcal{G}$ (also called ``group elements''), a base element $\ngen \in \mathcal{G}$, a prime $p$, a finite set of honest exponents $\honestexponent \subset \mathbb{Z}$, a finite set of exponents $\exponent \subset \mathbb{Z} / p\mathbb{Z}$, and an efficiently computable exponentiation function $\nexp : \mathcal{G} \times \mathbb{Z} \foutput \mathcal{G}$, where we write $X^y$ for $\nexp(X, y)$.
        The exponentiation function is required to have the following properties:

        \begin{enumerate}
        \item $(X^y)^z = X^{yz} $ for all $ X \in \mathcal{G}, y, z \in \mathbb{Z}$.
        \item The function $\phi$ defined by $\phi(x) = g^x$ is a bijection from $\exponent$ to $\{\ngen^x | x \in [1, p - 1]\}$.
        \end{enumerate}
    \end{definition}

    As done in \cite{EC:ABHKLR21} for a nominal group $\nominalgroup = \nominalgroupexpanded$, we define $D_H$ to be the uniform distribution of honestly generated exponents $\honestexponent$ and $D_U$ to be the uniform distribution on $\exponent$.
    We also recall the definition for the statistical distance between these distributions: $\Delta_{\nominalgroup} = \Delta [ D_H, D_U ] = \frac{1}{2} \underset{x \in \mathbb{Z}}{\sum} | \underset{D_U}{Pr}(x) - \underset{D_H}{Pr}(x) | $ \cite{EC:ABHKLR21}.

    \subsection{Strong Diffie-Hellman Problem}
    \label{sec:strong-dh}

    \begin{definition}[Strong Diffie-Hellman (\sdh) Problem]
        We define the advantage function of an adversary $\adversary{A}$ against the Strong Diffie-Hellman problem over the nominal group $\nominalgroup$ as

        \[
            \advantage{\nominalgroup, \adversary{A}}{\sdh} = \Pr \left [Z = g^{xy} \left | x,y \assign \varepsilon_u; Z \assign \adversary{A}^{DH(\cdot, \cdot)}(g^x, g^y) \right. \right].
        \]
        %
        where DH is a decision oracle that on input $(Y, Z)$ with $Y, Z \in \mathcal{G}$, returns $1$ iff $Y^x = Z$ and $0$ otherwise.
    \end{definition}


    \section{Introducing HQB}\label{sec:hqb}

    We consider a construction which combines a $\KEM$ and a nominal group to form a hybrid $\KEM$.


    \begin{definition}[\hqb] \label{def:hqb}

        Let $\nominalgroup = \nominalgroupexpanded$ be a nominal group, let $\kem$ be a key encapsulation mechanism and let $H$ be a hash function.
        We define the \hqb \kem as depicted in \autoref{fig:hqb}.

        \begin{gamefigure}{\hqb Framework.}{fig:hqb}
            \begin{games}{2}
                \begin{gameenv}{\algor $\keygen()$}
                    \State $(\sk_1, \pk_1)  \sample \kem.\keygen()$
                    \State $\sk_2 \sample \honestexponent$
                    \State $\pk_2 \assign \nexp(\ngen, \sk_2)$
                    \State $\pk \assign (\pk_1, \pk_2)$
                    \State $\sk \assign (\sk_1, \sk_2, \pk)$
                    \State \Return $(\sk, \pk)$
                \end{gameenv}
                \columnbreak
                \begin{gameenv}{\algor $\enc(\pk)$}
                    \State $(\pk_1, \pk_2) \assign \pk$
                    \State $k_1, c_1 \sample \kem.\enc(\pk_1)$
                    \State $\sk_e \sample \honestexponent$
                    \State $c_2 \assign \nexp(g, \sk_e)$
                    \State $k_2 \assign \nexp(\pk_2, \sk_e)$
                    \State $c \assign (c_1, c_2)$
                    \State $k \assign H(\text{label} \concat k_1 \concat k_2 \concat c \concat \pk)$
                    \State \Return $(k, c)$
                \end{gameenv}
            \end{games}

            \begin{games}{2}
                \begin{gameenv}{\algor $\dec(c, \sk)$}
                    \State $(\sk_1, \sk_2, \pk) \assign \sk$
                    \State $(c_1, c_2) \assign c$
                    \State $k_1 \assign \kem.\dec(c_1, \sk_1)$
                    \State $k_2 \assign \nexp(c_2, \sk_2)$
                    \State $\ifthen{k_1 \compare \bot}$
                        \State \quad \Return $\bot$
                        \State $k \assign H(\text{label} \concat k_1 \concat k_2 \concat c \concat \pk)$
                        \State \Return $k$
                    \end{gameenv}
                    \gamebreak
                    \emptygame
                \end{games}
            \end{gamefigure}

        \end{definition}

        \subsection{Correctness}

        The correctness of the HQB framework follows from the correctness of the \kem used to instantiate \hqb.
        Formally, if \hqb is instantiated with a $\delta$-correct \kem, then the correctness of \hqb is defined as follows:

        \[
            \Pr \left[\,  k \ne k' \, \left | \substack{(\sk, \pk) \sample \keygen(\,) \\ (k, c) \sample \enc(\pk) \\ k' \gets \dec(c, \sk) } \right. \right] \le \delta .
        \]


        \section{Security of HQB}\label{sec:hqb-proofs}

        The objective of a hybrid of this form is that the hybrid remains secure so
        long as one of the component algorithms remains secure. In other words, one
        algorithm may be found to contain weaknesses without compromising the hybrid.

        \paragraph{Proof-strategy.}

        Our proofs cover multiple scenarios, including the possibility of
        cryptographically-relevant quantum computers becoming a reality, and thus
        breaking the $\indcca$ security of pre-quantum KEMs, and the possibility of
        post-quantum cryptography such as a KEM being classically broken, such as by
        an attack, vulnerability, or implementation error.

        \medskip \noindent
        \textsc{Pre-Quantum Security.} This the case where one of the component
        KEMs is assumed to be classically $\indcca$-secure but where the post-quantum
        KEM may not offer any $\indcca$ security.

        We prove in the standard model that HQB is $\indcca$-secure when:
        \begin{enumerate}
        \item The KDF is modeled as a secure $\PRF$;
        \item The pre-quantum component is $\sdh$-secure; and
        \item The post-quantum KEM is modeled as a possibly broken KEM that provides no security.
        \end{enumerate}

        Our proof also demonstrates the opposite scenario:

        \medskip \noindent
        \textsc{Post-Quantum Security:} This the case where the component
        KEMs is assumed to be $\indcca$-secure against a quantum adversary but the
        pre-quantum component group may not offer any security at all.

        We prove in the standard model that HQB is $\indcca$-secure when:
        \begin{enumerate}
        \item The KDF is modeled as a $\prf$;
        \item The pre-quantum component group no longer provides any security; and
        \item The post-quantum KEM provides $\INDCCA$ security against a quantum attacker.
        \end{enumerate}

        These results do not rely on the ROM nor the QROM: it is sufficient to assume
        the standard model $\PRF$ security property for the employed hash function
        $F$.

        \subsection{Reduction to $\sdh$ security of $\nominalgroup$}

        Starting assumptions:

        \begin{itemize}
        \item $\nominalgroup$ is $\sdh$-secure and
        \item $F$ is a secure split-key PRF
        \end{itemize}

        \todo[inline]{Literally the SDH/ROM model proof of X-Wing except we do
            include the whole hybrid pk and ct so we can basically ignore all the
            security properties of the broken KEM as long as we can model the PRF as
            a secure split-key PRF}

        \subsection{Reduction to $\indcca$ security of $\KEM$ and PRF security in the standard model}

        \todo[inline]{Literally the standard model proof of X-Wing relying on the
            IND-CCA security of the KEM and the KDF being a secure PRF}


        \section{Concrete Instantiations}\label{sec:concrete}

        \subsection{NAME: $\HQB[ML-KEM-768, P-256, HKDF-SHA256]$}

        \todo[inline]{specify}

        \subsection{NAME2: $\HQB[ML-KEM-768, X25519, SHA3]$}

        \todo[inline]{specify}

        \todo[inline]{
            Ref the IND-CCA proofs for ML-KEM from the FIPS submission or wherever.

            Ref the nominal group proofs for x25519 and the short weierstrass curves from several x-wing and starfighters

            Ref whatever we got for SHA3 as a nice keyed PRF (keccak docs?)

            Ref whatever we got for HKDF-SHA256 as a secure keyed PRF (Krawczyk?)
        }

        \bibliographystyle{plain}
        \bibliography{local,cryptobib/abbrev0,cryptobib/abbrev3,cryptobib/crypto}

    \end{document}
